\begin{abstract}
We characterize the optimal expected length of uniformly honest confidence intervals for a target-population complier average treatment effect in an equal-size two-sample experiment. The data comprise \(n\) independent source records containing a scalar covariate, encouragement, treatment receipt and outcome, together with \(n\) independent target covariate records. Source and target densities, the encouragement propensity and conditional arm means are unknown and belong to fixed Hölder classes with exponent \(1/8\). For \(n\ge256\), under fixed density envelopes, encouragement overlap, bounded potential outcomes, instrument validity, monotonicity, and transport of conditional complier shares and complier-weighted outcome effects, the minimax expected length is of order \(\min\{1,n^{-1/3}/a\}\), uniformly over \(0<a\le1/4\), where \(a\) is the transported first-stage evaluation floor. Coverage is required over the full model with positive transported compliance. An observable split-sample cubic estimator and calibrated affine-score inversion attain this order with a single interval sequence across all floors. A matching lower bound uses a continuum of admissible instrumental-variable laws with a no-defiers potential-outcome realization and applies to arbitrary measurable confidence sets. The frontier has a constant-order branch when the strength floor is at or below the rough estimation scale and an inverse-strength branch above it.
\end{abstract}

\section{Introduction}
\label{sec:introduction}

Transporting an instrumental-variable effect requires learning both the population weights and the encouragement contrasts that identify the effect among target-population compliers. When transported compliance is small, uncertainty in those contrasts is amplified by the compliance denominator. This paper characterizes that interaction through the optimal expected length of confidence intervals with finite-sample coverage uniformly over a specified continuous-covariate model.

Our experiment combines two independent samples of equal size: source records containing covariates, encouragement, receipt and outcomes, and target records containing covariates. The covariate is scalar and lies in the unit interval. Source and target densities, the source encouragement propensity, and outcome and receipt arm means belong to fixed Hölder classes with exponent \(1/8\). Density envelopes are fixed and positive, encouragement probabilities lie between \(1/4\) and \(3/4\), and potential outcomes lie in \([0,1]\). Conditional instrument randomization, consistency, exclusion and monotonicity connect observed source contrasts to complier quantities. Two transport equalities, holding for target-almost every covariate value, carry the conditional complier share and the conditional complier-weighted outcome effect across populations. Under these restrictions, \cref{obj:lem:transported-cace-identification} identifies the target complier average treatment effect as the ratio of target-averaged source outcome and receipt contrasts.

The main result concerns intervals whose coverage is at least \(1-\alpha\) at every law in this model, where \(\alpha\in(0,1)\) is the fixed noncoverage probability. Precision is evaluated on laws whose transported first stage lies between \(a\), the evaluation floor, and \(1/4\). Thus global coverage and strength-specific expected length are distinct requirements. For fixed density envelopes and Hölder radius, \cref{obj:thm:sharp-length-frontier-full-elbow} establishes
\[
L_n(a)\asymp \min\{1,n^{-1/3}/a\},
\qquad n\ge256,\quad 0<a\le1/4,
\]
where \(L_n(a)\) is the minimax worst-case expected length among globally honest connected intervals. The comparison constants depend only on the fixed regularity parameters and noncoverage probability. The bounds are uniform across the stated strength range and contain no leading logarithmic factor.

A single observable interval sequence attains this frontier. The construction expresses the transported contrasts as smooth rational functionals of seven observable marked densities. A clipped histogram pilot, followed by linear, quadratic and cubic corrections using separate evaluation blocks, yields the uniform mean-square bound in \cref{obj:lem:marked-cubic-mse}. Inverting an affine score with a deterministic radius calibrated from that bound gives the interval in \cref{obj:def:score-interval}. Its coverage holds for every positive sample size, and its expected length is bounded by a constant times the rough estimation scale divided by the actual transported first stage, capped at constant order. The evaluation floor enters the performance criterion; the same procedure serves every strength slice.

The converse preserves the causal structure of the problem. The construction in \cref{obj:def:legal-iv-mixture,obj:lem:legal-iv-mixture-full-elbow} supplies admissible observed probability laws together with compatible potential receipts and outcomes satisfying monotonicity and the transport restrictions. Its actual compliance share is the larger of the evaluation floor and the rough estimation scale. A continuum of common effect values across sign mixtures, combined with observational proximity to a center experiment, produces the matching expected-length lower bound. \Cref{obj:thm:honest-lower-full-elbow} applies to arbitrary measurable confidence sets contained in the bounded effect domain. The constant-order branch therefore reflects the information available under global coverage when the strength floor is at or below \(n^{-1/3}\).

This characterization joins three established lines of work. Transported complier-effect identification and regular inference are developed by \citet[Theorems 3.1, 4.2 and 4.3]{ChenHuang2025Generalizing}, while \citet[Theorem 6]{Clark2024} establish regular inference for transported principal effects under their identification and nuisance-rate conditions. Rough transport-functional estimation has a direct antecedent in \citet[Theorems 6 and 8]{ZengKennedyBodnarNaimi2025Efficient}; their lower-bound exponent specializes to one third under the stated one-dimensional smoothness and design conditions. Higher-order correction follows the functional-estimation approach of \citet[Section 1]{Robins2008}. Identification-robust score inversion is developed for LATE by \citet[Theorem 3.2]{Ma2026IdentificationRobust}, and its confidence-set geometry is studied by \citet[Section 3 and Theorems 1 and 2]{SmuclerLanniMasip2025Score}. The immediate expected-length benchmark is \citet{CausalSmith2026TransportedLATEStrengthFrontier}, which treats supplied geometry and specified cell-based weight-learning settings. The present result connects rough functional estimation and score inversion to a sharp expected-length criterion with fully learned continuous geometry, equal sample sizes and uniform coverage under the declared causal model.

The analysis covers scalar covariates, equal source and target sample sizes, fixed Hölder exponent (1/8), and the causal and transport restrictions stated in \cref{sec:model}. \Cref{sec:related-work} gives the closest methodological comparisons; \cref{sec:model} specifies the observed experiment and inference criterion; and \cref{sec:main-results} states the frontier and its mechanism. \Cref{sec:discussion} interprets the calibration and formulates the unequal-sample problem. The technical construction and supporting bounds appear in \cref{sec:upper-appendix}, the legal lower family in \cref{sec:lower-appendix}, and the auxiliary comparators and verification scope in \cref{sec:verification-appendix}. Full proofs of the anchored results are collected in \cref{sec:deferred-proofs}.
